% TOP_PROBLEM: {"id": 2754, "title": "Kirby Problem 2.6", "category_id": 7, "category": {"id": 7, "name": "topology", "display_name": "Topology", "description": "Properties preserved under continuous deformations.", "slug": "topology", "order_index": 7, "created_at": "2026-07-31T15:26:25.671Z"}, "status": "open", "problem_number": "KP-2.6", "set_id": 11}
% FIRST_POSTED: 2026-10-01
% ENTRY_KIND: statement_only
% UnsolvedMath ID 2754; source catalogue code KP-2.6
% !TeX program = lualatex
% Definition and sources only; this file is not an LLM solution attempt.
\documentclass[11pt,a4paper]{article}
\usepackage[margin=25mm]{geometry}
\usepackage{fontspec}
\setmainfont{lmroman10-regular.otf}[BoldFont=lmroman10-bold.otf,
 ItalicFont=lmroman10-italic.otf,BoldItalicFont=lmroman10-bolditalic.otf]
\usepackage{amsmath,amssymb,mathtools}
\usepackage{unicode-math}
\setmathfont{latinmodern-math.otf}
\AtBeginDocument{\let\setminus\smallsetminus}
\usepackage{xurl,hyperref}
\hypersetup{colorlinks=true,urlcolor=blue}
\setcounter{secnumdepth}{0}
\setlength{\parindent}{0pt}
\setlength{\parskip}{5pt}
\setlength{\emergencystretch}{3em}
\begin{document}
\section*{Kirby Problem 2.6}\label{2754}
{\small UnsolvedMath ID 2754 · KP-2.6 · Topology · Kirby's Problems in Low-Dimensional Topology}
\subsection{Definitions and mathematical statement}
For a compact oriented surface $S_g$ of genus $g$, let $\operatorname{Mod}(S_g)$
be the group of orientation-preserving homeomorphisms up to isotopy. If $c$ is an
essential simple closed curve, its Dehn twist $T_c$ is the mapping class obtained
by rotating an annular neighborhood of $c$ once and extending by the identity.
A multitwist is a product $\prod_{i=1}^r T_{c_i}^{n_i}$ along pairwise disjoint
curves, where the $n_i$ are integers. Powers of individual twists are included.

The problem is to classify and enumerate the finite-index subgroups of
$\operatorname{Mod}(S_g)$ that admit generating sets consisting of Dehn twists,
multitwists, or powers of these, distinguishing these generating-set restrictions.
One should specify the equivalence used for the classification, for example
conjugacy inside the mapping class group, as well as the genus and subgroup index.

A related part asks which such subgroups are stabilizers of naturally specified
geometric or topological data on the surface, and conversely which stabilizers
have the stated generation property. For an action on a space of structures and
a structure $q$, its stabilizer is $\{f:f\cdot q=q\}$. Examples of the intended
data include spin structures and homological data. The source's request for a
geometric classification is a research program, rather than a prescribed list of
allowable structures; an arbitrary coset action alone is not the intended explanation.
\subsection{Short English statement}
Classify finite-index mapping class subgroups generated by twists, and explain which arise by preserving natural structures on a surface.
\subsection{Sources}
\begin{itemize}
\item[S1] ULAM.AI and the UnsolvedMath Contributors, supplied problems.json, record 2754 (KP-2.6), Kirby's Problems in Low-Dimensional Topology. Catalogue identity and source transcription; CC BY 4.0.
\url{https://huggingface.co/datasets/ulamai/UnsolvedMath}.
\item[S2] K3: A New Problem List in Low-Dimensional Topology, Problem 2.6. Expanded formulation of the supplied catalogue statement.
\url{https://math.berkeley.edu/sites/default/files/surv-295-ruberman-watermarked-author-pdf.pdf}.
\end{itemize}
\end{document}
